\documentclass[11pt]{article}

\input{../shared/preamble.tex}

\usepackage{amsmath}
\usepackage{booktabs}
\usepackage{cleveref}

\title{Repaired Hierarchical Portfolios}
\author{Peter Cotton}
\date{August 25, 2026}

\newcommand{\E}{\mathbb{E}}
\newcommand{\Prob}{\mathbb{P}}
\newcommand{\R}{\mathbb{R}}
\DeclareMathOperator*{\argminop}{arg\,min}
\DeclareMathOperator*{\argmaxop}{arg\,max}

\newtheorem{proposition}{Proposition}
\newtheorem{remark}{Remark}

\begin{document}
\maketitle

\begin{abstract}
Hierarchical Risk Parity and its Schur-complementary generalization build a
portfolio by recursive bisection: only the covariance within the current
block is consulted at each split, and at Schur's coupling $\gamma=0$ (plain
HRP) cross-block covariance is never consulted at all. This is what makes
the construction robust and inversion-free, but it also means the weights
are, structurally, a function of less information than is available. We
repair this without re-optimizing: calibrate the latent Thurstone abilities
that reproduce the hierarchical weights under a reconstruction of the
dependence the recursion actually used, a nested factor model over the same
median-bisection tree, then re-race those scores under the fuller dependence
on hand. On a known-truth contagious market, correlation completion cuts
out-of-sample expected shortfall by $10$--$17\%$ relative to plain HRP and
Schur, survives sample sizes down to two observations per asset, and needs
less repair as Schur's own $\gamma$ rises. The same comparison on
$12{,}000$ real S\&P~500 and REIT trials confirms it: correlation completion
wins $59$--$80\%$ of trials across five covariance estimators, and on real
data the gain turns out to come mostly from the race's own reallocation, not
from correlation recovery per se. A companion test of whether genuine tail
dependence adds value beyond correlation completion found no evidence for
it on the synthetic market, and a real-data result that splits by universe
rather than confirming or reversing that null; we report both.
\end{abstract}

\section{Introduction}

Hierarchical Risk Parity \parencite{lopezdeprado2016} and the
Schur-complementary generalization \parencite{cotton2024schur} never invert a
covariance. That is their answer to mean--variance's central fragility,
inverting an ill-conditioned one. A seriation orders the assets so that
similar ones sit near each other. Recursive bisection then allocates between
the two halves of each split in inverse proportion to each half's own
variance, augmented by the Schur complement of the other half scaled by a
coupling parameter $\gamma\in[0,1]$: $\gamma=0$ ignores the other half
entirely (plain HRP), and $\gamma\to1$ tends to the minimum-variance
portfolio.

The robustness has a structural cost. At $\gamma=0$, the recursion never
reads the covariance \emph{between} the two halves of any split, not because
that covariance is estimated to be small, but because the construction does
not consult it. At $\gamma>0$ some of it is recovered through the Schur
complement, but never all of it short of $\gamma=1$, and $\gamma=1$
reintroduces the inversion HRP exists to avoid. Every hierarchical portfolio
short of full minimum variance is, by construction, built from
\emph{less} dependence information than is on hand.

This paper asks whether that discarded information can be added back without
re-optimizing, using the ability tilt of \textcite{cotton2026polish}. Its
central operation is built for exactly this: calibrate a latent score vector
that reproduces a given allocation under one dependence law, then
re-evaluate the same score vector under a different law. Read the
hierarchical weights as a benchmark, calibrate under a reconstruction of
\emph{the dependence the recursion actually consulted}, and re-run the race
under the fuller dependence on hand. If the reconstruction is reasonable, the
repair adds back specifically what was discarded, and nothing more.

The paper's contribution is fourfold. First, a constructive object for ``the
dependence a hierarchical bisection actually consulted'': a nested factor
model with one latent factor per retained internal split of the same
median-bisection tree HRP/Schur already build, so it drops directly into the
$k$-factor calibration engine of \textcite{cotton2026factorprobit} with no new
machinery. Second, a known-truth simulation (correlated pins that cascade in
physically clustered groups, producing genuine tail contagion a covariance
cannot see) on which the repair is tested against plain HRP and Schur, across
the amount of tree retained, across Schur's own $\gamma$, and across sample
sizes small enough that the covariance is honestly noisy rather than known.

Third, a replication on real markets: the same comparison across $12{,}000$
random trials on S\&P~500 and REIT sub-portfolios, varying every choice the
repair itself does not control (estimation window, covariance estimator,
sub-portfolio, $\gamma$). It confirms that correlation completion
generalizes, and it sharpens rather than reverses the synthetic-market
finding that most of the gain is a race-reallocation effect, not correlation
recovery per se. Fourth, an honest accounting of what worked and what did
not. Correlation completion helped substantially and survived every
robustness check we ran. A sharper test of whether genuine tail dependence
adds value \emph{beyond} correlation completion did not confirm the
hypothesis on the synthetic market, and on real data the same test splits
cleanly by universe (\Cref{sec:realmarket}) rather than confirming or
reversing the null outright; we report both rather than discard either.

\section{Background}

\paragraph{The ability tilt.} Following \textcite{cotton2026polish}, write
$\theta=-a$ for minus the latent ability vector, and
\[
  G_S(\theta) = \E_{\eta\sim S}\big[\max_i(\theta_i+\eta_i)\big], \qquad
  W_S(\theta) = \nabla G_S(\theta) \in \Delta,
\]
so $W_S(\theta)$ is the vector of winning probabilities of a race among
competitors with strengths $\theta$ under a centered dependence law $S$. For
a positive-definite Gaussian reference $S_0$, $W_{S_0}$ is a smooth bijection
from $\R^n/\mathrm{span}\{\mathbf 1\}$ onto the open simplex, so every
strictly positive benchmark $w^0$ has a unique calibrated score
$\theta^0=W_{S_0}^{-1}(w^0)$. Re-evaluating that same $\theta^0$ under a
different law $S_1$ gives the tilt $T_{S_0\to S_1}(w^0)=W_{S_1}(\theta^0)$,
which reproduces $w^0$ exactly when $S_1=S_0$.
\textcite{cotton2026factorprobit} supplies the computational engine this
paper uses throughout. For a $k$-factor Gaussian law
$X_i=\mu_i+v_i^\top f+\sqrt{D_i}\,\varepsilon_i$, both $W_S$ and $W_S^{-1}$
are computable in $O(QNL)$ (quadrature nodes $\times$ assets $\times$ lattice
points), with a damped Jacobi calibration and analytic own-location slopes
replacing the far slower fixed-point inversion this paper's experiments would
otherwise require.

\paragraph{Hierarchical and Schur-complementary allocation.} Given a
seriation order and a covariance $\Sigma$, the Schur-complementary
construction \parencite{cotton2024schur} recursively bisects the ordered
asset list. At each split into halves $L,R$, each half's own block
$\Sigma_{LL}$ (resp.\ $\Sigma_{RR}$) is augmented by the $\gamma$-scaled Schur
complement of the other half,
\[
  \Sigma_{LL}^{\mathrm{aug}} = \Sigma_{LL} - \gamma\,\Sigma_{LR}\Sigma_{RR}^{-1}\Sigma_{RL}
  \quad(\text{and symmetrically for } R),
\]
and the split allocates between $L$ and $R$ in inverse proportion to each
augmented block's naive (inverse-variance) portfolio variance. At $\gamma=0$
the augmentation is the identity map, $\Sigma_{LR}$ is never read, and the
recursion is exactly Hierarchical Risk Parity. Every off-diagonal entry of
$\Sigma$ is consulted only while its two assets remain in the same block. The
instant they are placed in different halves, their covariance is never read
again (at $\gamma=0$) or read only through the reduced Schur term (at
$\gamma>0$).

\section{Method}
\label{sec:method}

\subsection{The repair operator}

Let $w^{\mathrm{base}}$ be the weights of HRP or Schur-complementary
allocation, built from an estimated covariance $\hat\Sigma$ and seriation
order $\mathrm{ord}$. Let $C^{\mathrm{used}}$ be a reconstruction of the
dependence the recursion consulted (\Cref{sec:tree}), and let
$C^{\mathrm{full}}$ be the fuller dependence on hand: the same $\hat\Sigma$
used as an ordinary correlation matrix, or a genuinely tail-dependent
empirical law. The repair is
\[
  \theta = W_{C^{\mathrm{used}}}^{-1}\big(w^{\mathrm{base}}\big), \qquad
  w^{\mathrm{repair}} = W_{C^{\mathrm{full}}}(\theta).
\]
Operationally: calibrate the latent scores that reproduce the hierarchical
weights under what the hierarchy used, then re-race those same scores under
what is actually available. No optimization is re-run and no covariance is
inverted; the only inversion is the calibration step, which
\textcite{cotton2026factorprobit} supplies in closed, fast form for any
factor-Gaussian reference.

\begin{remark}[Fixed point]
If $C^{\mathrm{used}}=C^{\mathrm{full}}$ exactly, $w^{\mathrm{repair}}=w^{\mathrm{base}}$: the
repair is a controlled perturbation of the base allocation, not a
replacement for it. In practice $C^{\mathrm{used}}$ is only ever an
approximation of what the recursion effectively consulted
(\Cref{sec:tree}), so this holds in the limit of retaining the whole tree,
not at any finite truncation.
\end{remark}

\subsection{Reconstructing what the recursion used}
\label{sec:tree}

At $\gamma=0$, an off-diagonal entry $\hat\Sigma_{ij}$ is read by the
recursion at every ancestor split above the point where $i$ and $j$ are
finally separated into different halves, and never again after that. This is
exactly the defining property of a nested (ultrametric) covariance: one
latent shock per internal node of a tree, shared by every leaf in that
node's subtree, so that two leaves' covariance is the sum of the shocks
attached to their common ancestors. We construct this directly from the same
median-bisection tree HRP/Schur recurse over (splitting the seriated asset
list in half, repeatedly): at each internal node splitting members into
halves $L,R$, its own shock variance is the gain in the average $L$--$R$
cross-correlation not already explained by ancestors,
\[
  c_v^2 = \max\!\big(0,\ \overline{\mathrm{corr}(L,R)} - \text{(floor already explained by ancestors)}\big),
\]
and nodes are retained by \emph{gain} $c_v^2\,|L|\,|R|$ (the covariance mass
newly explained), highest first, up to a budget of \texttt{max\_factors}
retained splits. This gives, for any budget $F$, a valid factor model
$(V,D)$, with $V$'s $F$ columns the retained nodes' subtree indicators
scaled by $\sqrt{c_v^2}$. $F=0$ is the independent reference: no correlation
used at all, the naive ``polish from scratch'' baseline. Increasing $F$
restores progressively more of the tree, in order of how much covariance
each split explains, and the construction is guaranteed positive
semi-definite by being a sum of independent latent shocks, not an arbitrary
masking of $\hat\Sigma$.

Two things follow immediately from this construction and matter for reading
the results below. First, it plugs directly into the $k$-factor calibration
and forward engine with no modification: $F$ is exactly the factor rank $k$.
Second, it gives a single, principled dial for how much of the discarded
information to restore, which is the paper's answer to the question of what
$C^{\mathrm{used}}$ should be. \Cref{sec:discussion} discusses the several
alternative constructions this paper did not test.

\section{Experiments}

\subsection{A known-truth contagious market}

We use a fully-specified generative market already in this codebase: a
rigid-body bowling simulation read as a market (a pin is a firm; its
displacement under a thrown ball is a continuous return; pins are arranged in
randomized tight spatial clusters with gaps between them, so a throw that
misses leaves everyone quiet and a throw into a cluster cascades the whole
cluster through pin-pin collisions). This produces genuine, mechanically
generated tail contagion, whole-cluster co-moves that a Gaussian copula fit
to the same marginals and correlation underpredicts by one to four orders of
magnitude in the deep tail (already documented for this generator in this
codebase). It is a deliberately adversarial setting for any correlation-only
method.

For each random market we split the simulated path in half (in-sample /
out-of-sample), fit HRP and Schur-complementary ($\gamma=0.5$ unless swept)
on the in-sample covariance, repair as above, and evaluate out-of-sample
expected shortfall at the 5\% level (ES95, the mean of the worst 5\% of
realized portfolio returns; less negative is better). All reported numbers
average over multiple random markets (different cluster geometries and
seeds). Quadrature and Monte-Carlo resolutions are reduced from production
defaults for this exploratory sweep ($Q=9$ Gauss--Hermite nodes or a
$2^{10}$-point scrambled-Sobol rule beyond four factors, lattice
resolution 251, $2^{14}$ race paths); code is in \texttt{experiments/}, in
\texttt{schur\_thurstone\_repair.py}, \texttt{schur\_gamma\_sweep.py}, and
\texttt{estimation\_noise\_sweep.py}.

\begin{remark}[Losses, not returns, and a robust scale]
The tail-dependent repair race must be run on \emph{losses}, not returns
(the race convention is minimum-performance-wins), and on a \emph{robust}
standardization. Bowling returns are extremely leptokurtic: the typical
quiet-day move is $50$--$80\times$ smaller than the full-sample standard
deviation, which rare cascades inflate. Standardizing by the raw standard
deviation therefore crushes quiet days to near-zero noise, and mixing mean
centering with a median-based scale silently injects a large common offset
into every row. Both mistakes were made and caught in the course of this
work before the results below were produced. Both are instances of the
general rule that a target law must be matched in scale and location to the
reference before comparing tilts against it, not merely centered.
\end{remark}

\subsection{Correlation completion}
\label{sec:main-result}

\Cref{tab:main} sweeps the retained-tree budget $F$ against two repair
targets: \emph{gauss} re-races under the full in-sample correlation (no
tail dependence, isolating correlation completion alone); \emph{bowl}
re-races under the true in-sample bowling returns, bootstrap-resampled
(genuine tail dependence). Twenty random markets ($n=20$ assets, three
spatial clusters, cluster tightness swept).

\begin{table}[h]
\centering
\caption{Out-of-sample ES95 (less negative is better), mean over 20 markets.
Schur uses $\gamma=0.5$.}
\label{tab:main}
\footnotesize
\begin{tabular}{lccccccc}
\toprule
 & base & $F=0$ & $F=1$ & $F=2$ & $F=3$ & $F=5$ & $F=8$ \\
\midrule
HRP, gauss     & $-1.337$ & $-1.250$ & $-1.224$ & $-1.268$ & $-1.232$ & $\mathbf{-1.222}$ & $-1.225$ \\
HRP, bowl      & $-1.337$ & $-1.344$ & $\mathbf{-1.291}$ & $-1.339$ & $-1.365$ & $-1.352$ & $-1.351$ \\
Schur, gauss   & $-1.567$ & $-1.390$ & $-1.364$ & $\mathbf{-1.362}$ & $-1.400$ & $-1.413$ & $-1.424$ \\
Schur, bowl    & $-1.567$ & $-1.527$ & $-1.484$ & $\mathbf{-1.482}$ & $-1.517$ & $-1.539$ & $-1.551$ \\
\bottomrule
\end{tabular}
\end{table}

Correlation completion helps substantially and robustly: at its best budget,
HRP improves from $-1.337$ to $-1.222$ (8.6\%) and Schur from $-1.567$ to
$-1.362$ (13.1\%), both beating equal weight ($-1.260$) and approaching
minimum variance ($-1.164$) without ever inverting a covariance.

\begin{remark}[The gain splits into a global correlation effect and a
calibration effect]
A meaningful share of the improvement is present already at $F=0$, and this
is not a correlation-free effect. The re-race step
$w^{\mathrm{repair}}=W_{C^{\mathrm{full}}}(\theta)$ always uses the true,
full correlation, at every budget, $F=0$ included; only the
\emph{calibration} reference pretends independence at $F=0$. What $F=0$
isolates is that the race applies that correlation \emph{globally and all at
once}, every asset's weight computed jointly with every other's in one
evaluation, where HRP/Schur's recursion applies the same correlation only
\emph{locally and sequentially}, one split's naive-inverse-variance ratio at
a time, never revisiting a pair of assets once they land in different
halves. For HRP, $F=0$ alone recovers $76\%$ of the total gain to the best
budget; for Schur, $86\%$.

Raising $F$ adds a second, distinct source of improvement: telling the
\emph{calibration} step about correlation too, on top of the
always-correlation-aware re-race. The \emph{incremental} gain from that
(best budget minus $F=0$) is $0.028$ for both HRP and Schur, nearly
identical in absolute terms despite the very different $F=0$ shares. Both
sources of improvement are real and both use the true correlation. They
should not be conflated when the size of ``the correlation-repair result''
is quoted, and neither should be mistaken for a correlation-free effect.
\end{remark}

\paragraph{Tail dependence: a null result.} The \emph{bowl} row is
consistently \emph{worse} than \emph{gauss} at each method's own best budget
(HRP: $-1.291$ vs.\ $-1.222$; Schur: $-1.482$ vs.\ $-1.362$), the opposite of
the hypothesis that racing under the true contagious law, rather than merely
the completed correlation, would do better still. We do not think this
closes the question, for three reasons.

First, bowling returns move continuously with displacement on every
simulated day, not only when a barrier is crossed, so ordinary correlation
already captures much of the within-cluster co-movement here. The genuinely
tail-specific residual may simply be small in this generator, consistent
with this codebase's own earlier finding that a downside semicovariance
already recovers ``almost all'' of a comparable tail effect elsewhere,
leaving only a small residual that is not a second moment at all. Second,
ES95 with a few hundred out-of-sample rows inspects only the worst
$\sim\!5\%$ of them, which may not reach deep enough into the tail for
contagion to matter; ES99 or the probability of several names crashing
together are sharper instruments we did not yet run. Third,
bootstrap-resampling a modest in-sample window for the rare cascade rows is
itself noisy, and that variance cost may be outweighing a real but small
bias benefit. We record this as an open question, not a negative conclusion.

\subsection{Repair depth should track how much was already used}
\label{sec:gamma-result}

If the repair is genuinely filling in what a hierarchical method discarded,
the right amount of repair should \emph{shrink} as the method itself
retains more cross-block information: as Schur's $\gamma$ rises toward the
minimum-variance endpoint, there is nothing left to restore.
\Cref{tab:gamma} sweeps $\gamma\in\{0,0.25,0.5,0.75,1.0\}$ against the same
budget grid (gauss target only, 16 markets).

\begin{table}[h]
\centering
\caption{Out-of-sample ES95 by Schur coupling $\gamma$ and repair budget $F$; bold marks each row's best.}
\label{tab:gamma}
\footnotesize
\begin{tabular}{lccccccc}
\toprule
$\gamma$ & base & $F=0$ & $F=1$ & $F=2$ & $F=3$ & $F=5$ & best $F$ \\
\midrule
$0.00$ & $-1.325$ & $-1.261$ & $-1.231$ & $-1.258$ & $-1.221$ & $\mathbf{-1.209}$ & $5$ \\
$0.25$ & $-1.386$ & $-1.291$ & $-1.263$ & $-1.282$ & $-1.247$ & $\mathbf{-1.241}$ & $5$ \\
$0.50$ & $-1.590$ & $-1.405$ & $-1.379$ & $\mathbf{-1.364}$ & $-1.416$ & $-1.420$ & $2$ \\
$0.75$ & $-1.807$ & $-1.679$ & $-1.603$ & $\mathbf{-1.587}$ & $-1.689$ & $-1.650$ & $2$ \\
$1.00$ & $-1.728$ & $-1.704$ & $-1.686$ & $\mathbf{-1.665}$ & $-1.696$ & $-1.695$ & $2$ \\
\bottomrule
\end{tabular}
\end{table}

The best budget is non-increasing in $\gamma$ ($5,5,2,2,2$), exactly as the
repair-fills-a-gap account predicts: at low $\gamma$, Schur (like HRP)
consulted little cross-block covariance and rewards a deeper repair; at
$\gamma\ge0.5$ it has already recovered enough that only a shallow repair
helps and a deeper one starts double-counting the same information (compare
$F=2$ to $F=5$ at $\gamma=0.75$: $-1.587$ vs.\ $-1.650$, a real reversal).
This is evidence for treating $\gamma$ itself as informative about how much
repair a given Schur portfolio needs, rather than choosing $F$ the same way
regardless of $\gamma$, as \Cref{sec:main-result} otherwise does for
simplicity.

\subsection{Robustness to estimation error}
\label{sec:estnoise}

Every number above is built from an \emph{estimated} in-sample covariance,
never an oracle, but with a fairly generous window (500 observations for 20
assets). Since robustness to estimation error is HRP's own original
motivation, the repair needs to be checked specifically where that error is
large. \Cref{tab:estnoise} fixes a common out-of-sample window (the same 500
rows for every row of the table) and shrinks the in-sample window from 600
observations down to 40, two observations per asset, a genuinely adversarial
covariance estimate, at budgets $F\in\{1,2\}$ (12 markets).

\begin{table}[h]
\centering
\caption{Out-of-sample ES95 vs.\ in-sample window $T_{\mathrm{in}}$ (20 assets, fixed 500-row test window); ``gain'' is the best repair minus base.}
\label{tab:estnoise}
\small
\begin{tabular}{lccccc}
\toprule
$T_{\mathrm{in}}$ & $40$ & $80$ & $150$ & $300$ & $600$ \\
\midrule
HRP base            & $-1.536$ & $-1.370$ & $-1.337$ & $-1.346$ & $-1.331$ \\
HRP best repair      & $-1.330$ & $-1.228$ & $-1.181$ & $-1.208$ & $-1.196$ \\
HRP gain             & $0.207$  & $0.142$  & $0.156$  & $0.138$  & $0.135$ \\
\midrule
Schur base           & $-1.919$ & $-1.617$ & $-1.533$ & $-1.559$ & $-1.534$ \\
Schur best repair     & $-1.817$ & $-1.428$ & $-1.332$ & $-1.330$ & $-1.312$ \\
Schur gain           & $0.102$  & $0.189$  & $0.201$  & $0.230$  & $0.222$ \\
\bottomrule
\end{tabular}
\end{table}

The repair helps at every sample size tested, including the sparsest. For
HRP the gain is, if anything, largest exactly where the covariance is
noisiest ($0.207$ at $T_{\mathrm{in}}=40$ vs.\ $0.135$ at
$T_{\mathrm{in}}=600$): a cruder covariance estimate seems to make the
hierarchy's own naive split rule worse, leaving more for the
race-reallocation effect of \Cref{sec:main-result} to recover. Schur's gain
is smaller at the sparsest window and roughly doubles by
$T_{\mathrm{in}}=300$ before flattening, still positive throughout but not
simply monotonic. Neither pattern suggests the repair is fragile to
estimation noise; if anything HRP's suggests the opposite.

\subsection{Real-market validation}
\label{sec:realmarket}

Everything so far is a single synthetic generator. We repeat the correlation-
and history-repair comparison on two real return panels sourced from
\textcite{winningportstudy}: 583 S\&P~500 constituents (2014--2024) and the
49-name REIT sector (2000--2024). Each of 12{,}000 trials draws, independently,
a universe, a sub-portfolio of 12--40 names, a random in-sample window
(250--749 trading days) and a fixed 125-day out-of-sample window, a covariance
estimator (this package's own EWMA at two halflives, or \texttt{precise}'s
Ledoit--Wolf, OAS, or plain empirical covariance), and Schur's $\gamma$; $F$
is swept over $\{0,1,2,3\}$ within each trial. This varies every choice the
repair does not control, which is exactly the robustness question a single
hand-picked configuration cannot answer. All $192{,}000$ (trial $\times$
method $\times$ target $\times$ F) evaluations completed without error.

\begin{table}[h]
\centering
\caption{Real-market correlation completion (\emph{gauss} target), pooled
over both universes, 12{,}000 trials. Win rate is the fraction of trials with
positive out-of-sample ES95 gain.}
\label{tab:realmarket}
\small
\begin{tabular}{lcccc}
\toprule
 & $F=0$ & $F=1$ & $F=2$ & $F=3$ \\
\midrule
HRP, mean gain    & $\mathbf{0.00090}$ & $0.00041$ & $0.00036$ & $0.00037$ \\
HRP, win rate     & $\mathbf{78.8\%}$  & $69.5\%$  & $67.3\%$  & $68.6\%$  \\
Schur, mean gain  & $\mathbf{0.00041}$ & $0.00029$ & $0.00028$ & $0.00029$ \\
Schur, win rate   & $\mathbf{68.7\%}$  & $68.0\%$  & $66.5\%$  & $67.7\%$  \\
\bottomrule
\end{tabular}
\end{table}

Correlation completion replicates on real data, and the direction of
\Cref{sec:main-result}'s decomposition remark \emph{sharpens}: $F=0$, pure
race-reallocation with no correlation restored, is now the single best
budget for both methods, with a $78.8\%$ win rate for HRP. Real, noisy
correlation estimates apparently give the deeper tree budgets less to work
with than the clean synthetic covariance did. Restoring more of it does not
reliably add value once the correlation itself is estimated, only estimated
imprecisely. This is direct evidence that the race-reallocation effect, not
correlation completion, is the dominant and more robust component of the
gain in practice, and it argues for shipping $F=0$, or a small budget, as
the default rather than searching for a deeper one.

\emph{Estimator choice.} The gauss repair beats base for every one of the
five covariance estimators tried, with win rates from $59\%$ (empirical
covariance, Schur) to $80\%$ (OAS, HRP); the two shrinkage estimators
(Ledoit--Wolf, OAS) show both the highest win rates ($78$--$80\%$) and the
largest mean gains, consistent with \Cref{sec:estnoise}: a more heavily
shrunk (cruder, more diagonal-leaning) covariance estimate leaves the
hierarchy's own split rule worse off and the race-reallocation effect more
to recover.

\emph{Tail dependence, revisited.} Racing under the true bootstrap-resampled
in-sample history (\emph{hist}) is, pooled over both universes, roughly a
coin flip at $F=0$ ($55\%$ win rate for HRP, $54\%$ for Schur) and turns
mildly negative once any correlation is restored ($F\ge1$: $44$--$50\%$),
consistent with \Cref{sec:main-result}'s null result, not a reversal of it.
But pooling hides a real split. On the REIT sector, \emph{hist} beats base
$62$--$63\%$ of the time; on S\&P~500 sub-portfolios it beats base only
$33$--$39\%$ of the time, actively hurting on average.

A concentrated, sector-driven universe (real estate, exposed to a common
rate factor) apparently gives the bootstrap enough genuine, repeated
dependence structure to exploit. A broad, diversified sample of large-cap
names does not, and the bootstrap's variance cost (a modest in-sample window
standing in for the true tail law, per \Cref{sec:main-result}) dominates
there instead. The effect also decays with sub-portfolio size for HRP
($55\%$ win rate at 12--20 names, falling to $28\%$ at 30--40), the same
bootstrap-noise signature. Tail-dependence repair is therefore not
uniformly useless: it is specifically the broad, low-cluster-structure case
where it costs more than it recovers.

\emph{A real-data motivation for the $\gamma$-blend.} \Cref{sec:gamma-result}
found the best repair budget non-increasing in $\gamma$ on the synthetic
market. On real Schur portfolios (gauss target) the pattern is richer: at
$\gamma\in\{0,0.25,0.5\}$, $F=0$ is again best (mean gain $0.00084$,
$0.00099$, $0.00060$); but at $\gamma=0.75$ the mean gain at $F=0$ is
\emph{negative} ($-0.00007$), turning positive only from $F=2$ onward
($0.00004$, $0.00007$), and at $\gamma=1$ (minimum variance) $F=0$ is
sharply negative ($-0.00031$) while $F\in\{1,2,3\}$ are all mildly positive
($\approx0.0002$). A Schur portfolio built with substantial cross-block coupling has already
used most of the correlation. Repairing it against a reference that
discards all of it ($F=0$) makes it actively worse, not just suboptimal.
A shallow, non-zero budget matches what the portfolio actually used. This is exactly what the $\gamma$-blended reference
$C^{\mathrm{used}}(\gamma)=(1-\gamma)C^{\mathrm{tree}}+\gamma\,\hat\Sigma$
proposed in \Cref{sec:discussion} is for, and this result is real-data
evidence that it is needed, not only that it would be elegant.

\emph{Not an artifact of the seriation choice.} Every result above reads
$\mathrm{ord}$ from this package's own Schur-complementary estimator, which
orders assets by the Fiedler vector of the correlation-affinity graph rather
than \textcite{lopezdeprado2016}'s original single-linkage dendrogram (a
deliberate upgrade of this codebase, for smoother turnover under streaming
updates). Repeating the $F=0$ comparison with a classical single-linkage
quasi-diagonal order in place of the Fiedler one, on 2{,}000 further real
trials, gives a $69.3\%$ win rate (against $75.7\%$ for Fiedler on the same
trials); the per-trial gains correlate at $0.96$ across the two seriations
($88\%$ sign agreement), so it is largely the same trials benefiting under
either choice, not two unrelated effects that happen to average out
positive. The gap between the two widens at deeper budgets ($F=1,2$:
classical falls to a coin flip, $53\%$/$51\%$, while Fiedler holds at
$63\%$/$61\%$), plausibly because single-linkage's known chaining tendency
gives the nested tree noisier splits to read correlation off of at depth,
while $F=0$ does not lean on the tree's quality at all. The headline
finding is a property of the repair, not of this package's particular
seriation upgrade.

\emph{The gain is volatility reduction, not a return forecast.} Every
number reported so far is ES95, a tail-risk statistic. \Cref{tab:returnmetrics}
decomposes $1{,}500$ further real trials ($F=0$, gauss target) into mean
return, annualized volatility, Sharpe ratio, and ES95, base versus repair,
with the fraction of trials in which the repair wins each statistic.

\begin{table}[h]
\centering
\caption{Real-market $F=0$ repair, decomposed by statistic (1{,}500 trials).
Win rate is the fraction of trials in which the repair improves that
statistic over base.}
\label{tab:returnmetrics}
\footnotesize
\begin{tabular}{lcccc}
\toprule
 & mean return & volatility & Sharpe & ES95 \\
\midrule
HRP win rate   & $46.9\%$ & $82.5\%$ & $51.6\%$ & $79.9\%$ \\
Schur win rate & $46.9\%$ & $77.9\%$ & $50.6\%$ & $75.4\%$ \\
\bottomrule
\end{tabular}
\end{table}

Mean return is, if anything, a coin flip slightly \emph{below} even
($46.9\%$ for both methods), indistinguishable from noise rather than a
real effect in either direction. This is the construction working as
designed, not a disappointing result. We calibrate $\theta$ from
$w^{\mathrm{base}}$ alone and hold it fixed through the re-race; neither
HRP, Schur, nor the repair itself ever consults a mean-return estimate
anywhere in the pipeline. So the repair has no channel through which it
\emph{could} move expected return systematically, and a $47\%$ win rate is
what that looks like empirically.

Volatility, by contrast, is reduced in $78$--$83\%$ of trials: the one
channel the repair actually touches, since it is built entirely from the
dependence structure. Sharpe improves \emph{on average} (lower volatility
over an unchanged mean mechanically raises it) but only barely more often
than not ($51$--$52\%$) at the level of a single 125-day trial, since a flat
mean return still carries its own sampling noise that a modest volatility
reduction does not reliably overcome window by window. In one line: this is
a volatility- and tail-risk-reduction mechanism that costs nothing in
expected return, not a source of alpha, and the data support exactly that
reading.

\subsection{The $\gamma$-blended reference}
\label{sec:blend}

We built the $\gamma$-blended reference and tested it directly, rather than
leaving it as a proposal. $C^{\mathrm{tree}}$ is taken as the full retained
nested tree (\Cref{sec:tree} with no truncation, $F$ unbounded), giving a
dense $C^{\mathrm{used}}(\gamma)=(1-\gamma)C^{\mathrm{tree}}+\gamma\,\hat\Sigma$.
Because this is generally full rank, we factor-fit it (rank $8$, the
contrast-space heuristic of \Cref{sec:tree}'s engine) and calibrate through
the same lattice engine the $F$-budget repair already uses, so the
comparison is not confounded by calibration precision. We checked this
choice against a same-engine fixed-point test (calibrate and re-race under
the identical factor-approximated law): the residual is $3\times10^{-8}$ in
$\ell_1$, confirming the calibration itself is essentially exact and any
remaining gap is the covariance factor-fit's approximation error, not a
convergence failure.

\Cref{tab:blend} compares the $\gamma$-blend against $F=0$ on $60$ further
real trials per $\gamma$, engine-matched throughout.

\begin{table}[h]
\centering
\caption{Out-of-sample ES95 gain (repair minus base), $F=0$ versus the
$\gamma$-blend, $60$ real trials per $\gamma$.}
\label{tab:blend}
\footnotesize
\begin{tabular}{lccccc}
\toprule
$\gamma$ & $0.00$ & $0.25$ & $0.50$ & $0.75$ & $1.00$ \\
\midrule
$F=0$, mean gain      & $0.00077$ & $0.00096$ & $0.00052$ & $-0.00120$ & $-0.00065$ \\
$F=0$, win rate       & $66.7\%$  & $81.7\%$  & $61.7\%$  & $45.0\%$   & $45.0\%$   \\
blend, mean gain      & $0.00006$ & $0.00037$ & $0.00014$ & $-0.00003$ & $0.00004$  \\
blend, win rate       & $50.0\%$  & $70.0\%$  & $73.3\%$  & $53.3\%$   & $53.3\%$   \\
\bottomrule
\end{tabular}
\end{table}

The blend does what it was built for: at $\gamma=0.75$ it lifts a clearly
losing repair (mean gain $-0.00120$, $45\%$ win rate) to roughly break-even
($-0.00003$, $53.3\%$), and at $\gamma=1.0$ it flips a losing repair to a
winning one ($-0.00065\to+0.00004$, $45.0\%\to53.3\%$ win rate). This is not
a free upgrade. At low $\gamma$ the blend trades away magnitude: mean gain
falls to roughly a tenth of $F=0$'s at $\gamma=0$ and to under half at
$\gamma=0.25$, though win rate at $\gamma=0.5$ is actually higher under the
blend ($73.3\%$ vs.\ $61.7\%$). The honest reading is a trade-off, not a
strict improvement: the blend rescues the high-$\gamma$ failure at some cost
to $F=0$'s low-$\gamma$ edge.

That trade-off has an immediate consequence, since $\gamma$ is a known
quantity, chosen by whoever builds the Schur portfolio, not something that
has to be estimated. An \emph{adaptive} policy, $F=0$ for $\gamma<0.75$ and
the blend for $\gamma\ge0.75$, needs no new calibration: it is a
recombination of the same \Cref{tab:blend} numbers, using whichever row
applies at each $\gamma$. Averaged equally across the five $\gamma$ settings
tested, the adaptive policy reaches mean gain $0.00045$ and win rate
$63.3\%$, against $0.00008$ and $60.0\%$ for $F=0$ used everywhere and
$0.00012$ and $60.0\%$ for the blend used everywhere: roughly four to six
times the mean gain of either fixed policy alone. Since which repair
reference to use is a choice made \emph{before} seeing the out-of-sample
data (it depends only on $\gamma$, which is set at construction time), this
adaptive policy is available for free and dominates committing to either
$F=0$ or the blend alone.

\section{Discussion}
\label{sec:discussion}

\paragraph{Dependence repair, not risk repair.} The repair holds the
calibrated scores fixed and changes only the geometry they are re-raced
under. It inherits no guarantee of moving toward the minimum-variance
portfolio, still less of minimizing any particular risk functional. It
preserves what the base allocator already decided, which names deserve
credit, and only changes how that credit is expressed under a fuller
dependence picture. Where a Schur portfolio already sits at $\gamma=1$,
exactly reproducing minimum variance, there is by construction no
second-moment information left to restore: \Cref{tab:gamma}'s $\gamma=1.0$
row shows a shallow but still real residual gain, consistent with the
race-reallocation effect of \Cref{sec:main-result} rather than further
correlation completion.

\paragraph{Apply once; do not iterate.} The repair has no reason to be a
fixed point of itself. $w^{\mathrm{repair}} = w^{\mathrm{base}}$ only when
$C^{\mathrm{used}} = C^{\mathrm{full}}$ exactly (\Cref{sec:method}'s
fixed-point remark), and at $F=0$ that never holds. Feeding the $F=0$
repair's own output back in as a new $w^{\mathrm{base}}$ and repairing
again is correspondingly unstable. On 500 real trials, a single further
application still won $77.8\%$ of the time, close to the one-shot rate. But
$98\%$ of trials eventually drove a weight to exactly zero within a handful
of iterations, which the log-domain calibration cannot invert. We reran the
same trials at $32\times$ the Monte-Carlo resolution and got the same
collapse, so it is a genuine geometric effect, not a sampling artifact.

Here is why, traced through one case. Two real names are correlated at
$0.845$, with \emph{unequal} base weights ($0.0525$ and $0.0994$).
Red-bus/blue-bus \parencite{mcfadden1974} guarantees an even split only for
equal-strength duplicates (the companion paper's redundancy-consistency
proposition, \parencite{cotton2026polish}); these two names are not
equal-strength. For unequal strengths, a Gaussian race correctly sends the
weaker share toward zero as the correlation between them approaches one. A
single application, at the true correlation of $0.845$, cut the weaker
share by half: a real but bounded correction, not a collapse. Repeating the
same correction on its own output forgets that the true correlation is
fixed at $0.845$ and behaves as though it kept climbing every round,
because nothing anchors it back to the data. One application reads the
data; further applications read the previous application's opinion of the
data. That is why we run the repair once, as a single map from a base
allocation to a repaired one, and never iterate it.

\paragraph{What $C^{\mathrm{used}}$ should be.} This paper tests two
families. The default is the gain-pruned nested tree of \Cref{sec:tree},
indexed by a single truncation depth $F$, chosen for tractability: it drops
directly into the existing $k$-factor engine. For Schur specifically, the
$\gamma$-blended reference of \Cref{sec:blend},
$C^{\mathrm{used}}(\gamma)=(1-\gamma)C^{\mathrm{tree}}+\gamma\,\hat\Sigma$,
reuses Schur's own coupling parameter directly instead of a separately
chosen $F$, and \Cref{sec:blend} shows it fixes $F=0$'s failure at high
$\gamma$ at some cost to $F=0$'s edge at low $\gamma$, so an adaptive
policy that picks between the two using $\gamma$ alone beats committing to
either one. A flat $K$-cluster block-diagonal reference using the
\emph{exact} (not factor-approximated) within-cluster covariance remains
untested: it would be more faithful within each cluster, but requires a
separate factorization per cluster and does not share a rank budget as
cleanly across the whole universe.

\paragraph{Composing with further law changes.} Because the repair is an
instance of the general tilt $T_{S_0\to S_1}$, it inherits that
construction's composition properties. Repairing toward one law and then a
further, richer law gives the same endpoint as repairing directly to the
final law, and restricting the universe afterward commutes with the choice
of dependence law, provided the reference calibration is exact at each step
\parencite{cotton2026polish}. We rely on this only informally here, to note
that a repaired portfolio can be safely fed into a further Thurstone polish
(a market-factor tilt, a tail-copula overlay) without re-deriving
consistency from scratch, and leave a formal treatment specific to the
hierarchical setting to future work.

\section{Conclusion}

Hierarchical portfolios are robust exactly because they refuse to consult
all the available dependence information. That refusal need not be paid
for twice. The same latent scores that rationalize the hierarchical weights
under what the recursion used can be re-raced under what is actually on
hand, recovering a meaningful share of the gap to a fuller allocation
without inverting anything. The recovered amount tracks how much was
already used, falling as Schur's $\gamma$ rises on the synthetic market and
turning sharply negative at $F=0$ once $\gamma$ is large enough that
independence is the wrong reference on real data. Reusing $\gamma$ itself
through a blended reference fixes that failure, at some cost to the
independent reference's edge at low $\gamma$; since $\gamma$ is chosen
before the repair ever runs, an adaptive policy that picks between the two
using $\gamma$ alone gets both, at four to six times the mean gain of
committing to either one. The repair survives both the covariance being
honestly noisy rather than known, and $12{,}000$ random trials against real
S\&P~500 and REIT sub-portfolios varying every choice the repair does not
control.

On real data the dominant, most robust part of the gain turned out to be
the shallowest one, race-reallocation over correlation completion, which is
a sharper and more actionable finding than the synthetic market alone could
give. It applies once and should not be iterated: the same mechanism that
correctly de-weights a genuinely redundant pair has no stopping rule of its
own if fed its own output. Whether genuine tail dependence adds value on
top of ordinary correlation completion remains open in general, but it is
no longer a blanket null: it costs more than it recovers on a broad,
diversified universe and modestly pays off on a concentrated, sector-driven
one. We report both rather than pick the one that flatters the method.

\printbibliography

\end{document}
